\chapter{Argument Structure, Alignment, and Voice}\label{ch:argument-structure}

\section{The root and VP}
\label{sec:vm-root}

For a simple verb, the lexical stem realises a root with a verbal categoriser
$v$. Categorisation, argument introduction and case licensing are distinct
operations in this analysis: the categoriser supplies verbal category, while
the predicate's lexical frame and any argument-building heads determine its
participants. An actor-introducing layer belongs to the inner predicate,
above Appl where present and below Caus; its ordinary exponent is zero.
The outer orientation head consequently need not introduce the actor whose
case its configuration determines.
A finite verb requires Asp.Align, and a bare stem cannot constitute an
independent finite predicate. Bound stems also occur in licensed non-finite
and nominal constructions, whose selecting heads need not include finite
Asp.Align (§\,\ref{sec:nominalised}).

A preverb--root combination is analysed here as a derived lexical predicate
within this innermost domain. Its internal prefixal order does not introduce
another member of the suffixal functional sequence in
Table~\ref{tab:vm-projections}; the ordinary Voice, causative, and aspectual
heads combine with the resulting stem. Prefixation preserves the base's
thematic class and construction-selected stem grades: \ird{nočk-},
\ird{stoj-}, and \ird{vad-} take ordinary progressive \ird{-ime} in
\ird{lenóčkime}, \ird{rustójime}, and \ird{navádime}, respectively.
No additional aspectual exponent or secondary imperfectivisation intervenes.

The \synhead{VP} formed by the root + its arguments constitutes the innermost
layer of the verbal clause. In a transitive clause the internal argument
(patient) is merged as the complement of V:

\begin{center}
\begin{forest}
  syntax tree
  [\synhead{VP}
    [{\synhead{DP}\\{\footnotesize patient}}]
    [V
      [{\textit{root}}]
    ]
  ]
\end{forest}
\end{center}

\noindent Because \synhead{VP} is head-final
(§\,\ref{sec:th-head-finality}), this complement surfaces to the left of the
verbal head at every embedded level. The preverbal position of the object is
therefore available in the base-generated configuration of the
head-final \synhead{VP}. This derives patient-before-verb order, not
patient-before-actor order. A patient may originate in this complement position
and participate in a separate clause-internal dependency without acquiring an
aboutness-topic interpretation.


\section{Alignment domains and case selection}
\label{sec:vm-asp-case}

In finite indicative clauses, Asp.Align encodes viewpoint aspect and the
alignment domain. Together with valency and orientation morphology, it
determines the clause's \term{absolutive pivot}, the argument that bears
direct case. Its frequent interpretation as an
aboutness topic belongs to the independent information-structural system
(Chapter~\ref{ch:infstr}). Aspect and alignment are jointly realised in the obligatory
suffix, with case valuation active in the finite domain. Participant
nominalisation preserves the aspectual and participant-selection features
without finite case valuation (§\,\ref{sec:nominalised}).

The ordinary active frames distinguish two domains. In the ERG-aspect domain
(perfective and retrospective), the patient takes direct case and the agent
transitive case; ordinary transitives require \ird{-n-}, whereas an
inherently intransitive predicate or the intransitive reading of an
ambitransitive takes the bare aspect suffix. In the NOM-aspect domain
(continuous, progressive and contemplative), the agent takes direct case and
the patient indirect case. The patient remains a core argument in either frame.

Valency morphology modifies the argument inventory before orientation and case
are determined. A simple antipassive licenses only its actor as an independent
core participant, and that actor takes direct case. A simple applicative
selects the applied participant as direct; its actor takes indirect case in the
NOM domain and transitive case in the ERG domain, while any retained theme is
indirect. A higher causative establishes a new orientation domain, in which the
causer is direct in NOM aspects; in ERG aspects the selected participant is the
retained lexical theme, otherwise the applied participant, otherwise the
causee (§\,\ref{sec:vm-caus}). These are constructional extensions of the
ordinary split, rather than additional meanings of viewpoint aspect.

The alignment split is an established property of Iridian, whose aspectual
 viewpoint and case assignment are jointly realised in Asp.Align. Perfective
 and retrospective transitives with \ird{-n-} assign direct case to the patient
 and transitive case to the agent, while the imperfective active frame assigns
 direct case to the agent and indirect case to the patient.
 These orientations concern grammatical realisation: a bounded event does not
 semantically entail a patient topic, and an ongoing activity does not require
 an agent topic. The grammar therefore treats their discourse association as
 a default rather than deriving aboutness from event culmination.

Event construal and information structure consequently make distinct
 contributions: an event may be presented as completed while the utterance
 continues a discussion of its agent, just as an ongoing event may be
 introduced without establishing any participant as an aboutness topic.

The ordinary active frames illustrate this contrast.
In the NOM-aspect domain, the agent or sole argument bears direct case,
while the patient of an ordinary transitive appears in
the indirect. This is a nominative-accusative frame, with the direct case
marking the privileged participant and the indirect marking the core patient:

\pex
\a
\begingl
\gla Marek zemácime.//
\glb Marek.\Dir{} run-\Prog{}//
\glft \trsl{Marek is running.} \lingcontext{intransitive}//
\endgl
\a
\begingl
\gla Marek travě pjaštíme.//
\glb Marek.\Dir{} bread.\Ind{} eat-\Prog{}//
\glft \trsl{Marek is eating bread.} \lingcontext{transitive; P = \Ind{}}//
\endgl
\a
\begingl
\gla Nužo mužna.//
\glb city.\Dir{}.\Def{} be.beautiful-\Cont{}//
\glft \trsl{The city is beautiful.} \lingcontext{stative continuous}//
\endgl
\a
\begingl
\gla Nužo mužnélime.//
\glb city.\Dir{}.\Def{} be.beautiful-\textsc{inch}-\Prog{}//
\glft \trsl{The city is becoming beautiful.} \lingcontext{dynamic progressive}//
\endgl
\xe

In an ordinary ERG-domain transitive, by contrast, the patient is
placed in the direct case, and the agent takes the transitive case. This is an
ergative-absolutive frame; what Iridian grammar calls the \emph{transitive case}
is, in typological terms, the ergative, and the \emph{direct case} serves as
the absolutive. An ordinary transitive in this domain requires the orientation head
\ird{-n-} to select its patient; when antipassive
\ird{-aš-} is present instead, the agent takes the absolutive and the patient is
removed from core argument structure (§\,\ref{sec:vm-antipassive}):

\pex
\a
\begingl
\gla Travó Marku pjaštník.//
\glb bread.\Dir{}.\Def{} Marek.\Trs{} eat-\textsc{trs}-\Pf{}//
\glft \trsl{Marek ate the bread.} \lingcontext{ergative; P = \Dir{}}//
\endgl
\a
\begingl
\gla Marek pjášček.//
\glb Marek.\Dir{} eat-\Antip{}-\Pf{}//
\glft \trsl{Marek ate.} \lingcontext{antipassive; A = \Dir{}}//
\endgl
\xe

The ergative transitive licenses a core patient in direct case, whether that
patient is overtly expressed or recovered from context. In a narrative about
that participant, the clause can describe what happened to it, as in the
bread example above; the case frame also remains available when the agent is
independently topicalised (§\,\ref{sec:tc-marked-topics}). Suffixal anchoring
identifies an accessible referent without being required by the construction.
Neither direct case nor transitive Voice entails that the patient undergoes
a change of state: the construction equally accommodates predicates of buying
and perception.

The antipassive presents the agent's activity without independently licensing
a patient, although the lexical event may still involve a conceptual theme.
This frame suits activity-oriented contexts in which no referential patient is
asserted or recovered. It excludes an indirect-case patient as well as a
direct-case one (§\,\ref{sec:vm-antipassive-desc}); an ordinary transitive
with a contextually recovered patient retains a different argument structure.

The ergative transitive and antipassive differ in the realisation of the
 patient and the selection of the absolutive pivot. The former licenses a
 core patient; the latter removes the patient from core argument structure.
 Their discourse uses follow from those valency differences, while aboutness
 remains independently available to suitable referents.

Iridian's alignment is \term{aspect-directed}: ordinary clauses distinguish
patient-direct ERG frames from agent-direct NOM frames. Applicative and
patient-selecting constructions establish their own orientation within those
domains, so the aspect value is one condition on the case frame rather than a
complete specification of its pivot.

The transitive case of an ordinary ERG-domain agent is analysed here as
dependent on a distinct participant licensed in direct case. Applicatives and
causatives extend this relation to their constructionally selected participant,
which need not be the lexical patient. This dependent-case analysis describes
the Iridian mapping; the surface contrast alone does not determine a universal
algorithm for case assignment.

Aspect selection establishes a case frame and temporal viewpoint, whose
 interaction can favour different discourse continuations. The pivot is the
 default aboutness candidate where the context supports ordinary topical
 predication, but a speaker can continue an agent topic in a perfective
 transitive without changing the event's aspect. The linked-topic dependency
 retains the agent's transitive case and the patient's direct case
 (§\,\ref{sec:tc-marked-topics}).

The suffix \ird{-n-} realises the higher orientation head, outside causative
morphology where present and taking \ird{-an-} after a rhotic causative.
It selects the patient of an ordinary transitive or
the retained theme of a derived causative, otherwise its applied participant,
and otherwise its causee. This priority is not equivalent to selecting the
structurally closest non-causer (§\,\ref{sec:vm-caus}). In the ERG domain it
combines with the ordinary transitive frame; in the NOM domain its
patient-oriented use licenses a passive, with the agent restricted to optional
indirect expression (§\,\ref{sec:vm-passive}). Simple applicatives establish
their own orientation through \ird{-éb-}, without a second overt \ird{-n-}.

Antipassive morphology builds a reduced predicate before case is assigned.
Its occurrence is licensed by the completed ERG-aspect configuration; it does
not inspect a patient whose case has already been valued. A following
causative may introduce a new causer and an outer \ird{-an-}, although the
lexical patient remains suppressed (§\,\ref{sec:vm-antipassive}).

\section{Argument inventory, selection and case domains}
\label{sec:vm-selection-domains}

The direct-case participant, the dependency licensing the external argument's
case, and the position in which an argument is pronounced are distinct parts
of the analysis. The following notation makes the first two explicit without
presupposing a particular landing site. $A$ is the lexical actor, $T$ a
retained transitive theme, $B$ an applied participant, $C$ a causer, and $E$
the causee. In a causative, the base actor or sole participant becomes $E$;
a theme-like sole argument of an intransitive is not a separate $T$.

Argument structure is built before the completed construction determines its
orientation. V takes $T$ as its complement where the lexical frame licenses
one. Appl introduces $B$ in relation to the lexical event, below the layer
introducing $A$; Caus then introduces $C$ above that event predicate.
Antipassive instead removes the independently licensed lexical theme from the
base inventory. The ordering of argument-building and orientation heads is
\[
 \mathrm{VP}\ \subset\ \mathrm{ApplP}/\mathrm{VoiceP}_{\mathrm{AP}}
 \ \subset\ \text{inner actor layer}\ \subset\ \mathrm{CausP}\ \subset\
 \mathrm{VoiceP}_{\mathrm{orient}}\ \subset\ \mathrm{AspP},
\]
with absent derivational layers omitted. Appl and antipassive are alternatives
under the existing co-occurrence restriction, not two freely stackable heads.
The inner actor-introducing layer lies above Appl and below Caus when those
heads are present. Outer orientation is thus above Caus, although the selected
participant can originate deep inside Caus's complement.

The case domain is the completed predicate licensed by finite aspect,
orientation and Mood. In a causative this includes the outer causer and the
licensed embedded participants; the base does not independently assign a
second finite direct case. A simple indicative applicative selects $B$;
active NOM causatives select $C$; patient-oriented and ERG causatives apply
$T>B>E$. Directive Mood selects a retained theme in either reduced aspectual
domain, overriding simple indicative applicative selection
(§\,\ref{sec:mood-directive-valency}). These are role-sensitive restrictions
on the completed configuration, not consequences of reading its suffixes.

\begin{table}[htbp]
\footnotesize
\caption{Participant inventories and case obligations. Parentheses mark an
optional passive dependent, not an omitted core argument.}
\label{tab:pivot-case-obligations}
\begin{tabularx}{\textwidth}{L L L}
\toprule
Construction & Licensed inventory & Finite case obligations \\
\midrule
Ordinary ERG transitive & $A,T$ & $T$: direct; $A$: transitive \\
Active NOM transitive & $A,T$ & $A$: direct; $T$: indirect \\
Simple indicative applicative & $A,B$; $T$ if retained & $B$: direct; $A$: indirect in NOM, transitive in ERG; retained $T$: indirect \\
ERG causative of intransitive & $C,E$ & $E$: direct; $C$: transitive \\
ERG causative of transitive & $C,E,T$ & $T$: direct; $C$: transitive; $E$: indirect \\
ERG causative of applied transitive & $C,E,B,T$ & $T$: direct; $C$: transitive; $E,B$: indirect \\
ERG causative of applied intransitive & $C,E,B$ & $B$: direct; $C$: transitive; $E$: indirect \\
ERG causative of antipassive & $C,E$ & $E$: direct; $C$: transitive; no licensed $T$ \\
NOM patient-oriented passive & $T,(A)$ & $T$: direct; expressed $A$: optional indirect dependent \\
Directive with retained theme & $A,T$; optionally $B$, or $C,E$ in a causative & $T$: direct; overt actor or causer: transitive; $B,E$: indirect where present \\
\bottomrule
\end{tabularx}
\end{table}

The ERG dependency relates the highest external actor or causer to the
construction's distinct selected participant. Residual indirect arguments are
licensed in their own roles; they are not automatically additional competitors
for transitive case. Otherwise a causee above a retained theme would wrongly
receive the same transitive marking as the causer. NOM passive agent licensing
and active NOM patient licensing remain distinct despite their shared indirect
exponent. An unpronounced but contextually licensed core argument remains in
the inventory; antipassive suppression does not.

This specifies the required mapping rather than settling its computational
implementation. A configurational implementation must delimit the competing
nominals and specify the licensing occurrence and evaluation cycle. Case may
be evaluated before displacement or with reference to a lower chain occurrence;
pronounced order alone cannot determine it; compare Baker's dependent-case framework and Baker and
Vinokurova's distinction between dependent case and functional-head licensing
\cite{baker2015,bakervinokurova2010}. A successive-Agree implementation must
independently define the probe's interaction and stopping conditions over the
same inventory; Clem and Deal's Shawi analysis provides a comparative model
\cite{clemdeal2026}. Neither
is supplied by the statement that Asp.Align conditions alignment, and neither
follows merely from the presence or absence of overt person agreement.
The comparison must also establish when the embedded predicate remains
accessible to the outer computation; a completed finite case domain is not
by definition a single transparent phase.

Neutral ordinary finite order places the selected participant before the
remaining arguments and the verb last. This is an output requirement for a
derivation, not an additional case-assignment step or a universal order for
every directive and reduced construction. In particular, the existing overt
contrastive directive subject can precede its direct-case theme
(§\,\ref{sec:mood-directive-case}). Allomorphy and Vocabulary Insertion
realise a configuration already licensed in syntax; no case ending or
Asp.Align exponent can retroactively choose the participant to be displaced.

\section{Valency operations}

Three derivational operations may be applied to the verb stem before the
aspect-alignment suffix: the antipassive, the applicative, and the causative.
Where more than one is present, the order is fixed: antipassive or applicative
precedes the causative.

\subsection{The antipassive (\ird{-aš-})}
\label{sec:vm-antipassive-desc}\label{sec:vm-antipassive}

The antipassive \ird{-aš-} forms an agent-oriented activity predicate from a
transitive base without licensing an independent patient DP. In the perfective
and retrospective, its agent receives direct case; the patient remains part of
the conceptual interpretation of the activity but cannot be expressed within
the finite clause. Iridian restricts this construction to the ERG-aspect domain.
A transitive base in that domain consequently requires an overt derivation:
\ird{-n-} retains the patient, \ird{-aš-} suppresses its independent
realisation, and \ird{-éb-} selects an applied participant. Bare perfectives
belong to independently intransitive readings (§\,\ref{sec:ambitransitive}).

The traditional notation \ird{-aš-} identifies this Voice construction;
it does not require an invariant string \ird{aš} to be concatenated with
every stem. Voice is realised through the replacement or augmentation specified
in Table~\ref{tab:verbclasses}. Thus \ird{pjašt-} selects \ird{pjášč-},
whereas the additive classes retain their thematic consonant before the
sibilant. These are morphological exponents of antipassive Voice, independently
of PAL and of the velar selector used by other constructions. Any accompanying
root-vowel quantity alternation, such as \ird{vad-}/\ird{váž-}, is specified
with its vocalic grade rather than inferred from the consonantal replacement.
The stem \ird{pjášč-} combines with perfective \ird{-ek}:

\pex
\a
\begingl
\gla Marek pjášček.//
\glb Marek.\Dir{} eat-\Antip{}-\Pf{}//
\glft \trsl{Marek ate.}//
\endgl
\a
\begingl
\gla Marek houlam se pjášček.//
\glb Marek.\Dir{} spoon.\Ind{} \textsc{inst} eat-\Antip{}-\Pf{}//
\glft \trsl{Marek ate with a spoon.}//
\endgl
\a
\begingl
\gla *Marek travě pjášček.//
\glb Marek.\Dir{} bread.\Ind{} eat-\Antip{}-\Pf{}//
\glft \trsl{(intended: Marek ate bread.)}//
\endgl
\a
\begingl
\gla *Marek pjašték.//
\glb Marek.\Dir{} eat-\Pf{}//
\glft \trsl{(intended: Marek ate.)}//
\endgl
\xe

The instrumental phrase in the second example modifies the activity without
restoring a patient. The third example fails because indirect case cannot
license a patient excluded by the antipassive; the fourth lacks the overt
morphology required by the transitive base. A nominal dependent may separately
restrict an antipassive event under participant nominalisation, where a nominal
relational head licenses it (§\,\ref{sec:nominalised}). This nominal relation
does not extend to finite antipassive clauses.


In the syntactic analysis, antipassive formation combines the transitive event predicate with an inner
Voice head that closes the patient role without projecting an independent
patient DP. For \ird{pjašt-}, its interpretation can be represented as
\[
 \lambda x.\lambda e\,[\operatorname{eat}(e)\land
 \operatorname{Agent}(e,x)\land
 \exists y\,\operatorname{Theme}(e,y)].
\]
The existentially bound variable supplies the conceptual participant of eating;
it is unavailable for independent nominal realisation in the finite clause.
The representation is an analysis of Iridian's strict patient-suppression rule.
The availability of an oblique patient varies across antipassive constructions
cross-linguistically, as does their association with ergativity
\cite{heaton2020}; neither the suppression rule nor its aspectual restriction
follows from the term \term{antipassive} alone.

Asp.Align licenses inner antipassive Voice only in the ERG-aspect domain.
This dependency holds of the completed configuration, in which the higher Asp
head selects a compatible lower predicate. It does not require the lower head
to inspect a patient already assigned case by a not-yet-merged Asp head. The
remaining agent receives direct case through the resulting intransitive frame.
Embedded Asp imposes the same selectional restriction under participant
nominalisation, although finite case valuation is inactive there.

Antipassive Voice and Appl are separate heads, whose co-occurrence Iridian
excludes within one derived predicate. This is a language-specific selectional
restriction, rather than a consequence of two exponents occupying one universal
Voice position. Cross-linguistic combinations of antipassives and applicatives
are discussed by Basilico \cite[ch.~5]{basilico2026}. A higher causative may nevertheless select an antipassive
predicate (§\,\ref{sec:vm-caus}); its new causer does not restore the
suppressed lexical patient.

The antipassive presents an activity without an independently licensed patient,
which can support discourse centred on the actor or on the activity itself.
An ordinary transitive with an omitted patient differs because its patient
remains a referential argument recoverable from context. Choosing between these
constructions therefore concerns argument structure as well as discourse.

\subsection{The applicative (\ird{-éb-})}\label{sec:vm-appl}

The applicative \ird{-éb-} licenses one additional participant in relation to
the event and, in a simple indicative applicative, selects that participant as
the direct-case pivot. Its productive default is benefactive: the event is
performed for an individual's benefit, without entailing that the individual
performs the lexical action or receives its theme. The actor receives indirect
case in the NOM-aspect domain and transitive case in the ERG-aspect domain;
a retained lexical theme receives indirect case in either domain.

The following constructed contrast retains the same actor, theme, and
beneficiary. In the first clause, \ird{my} heads a benefactive phrase; in the
second, the beneficiary is a core DP selected by \ird{-éb-}:

\pex\label{ex:val-appl-transitive}
\a
\begingl
\gla Marek travě bylam my pjaštíme.//
\glb Marek.\Dir{} bread.\Ind{} child.\Ind{} for eat-\Prog{}//
\glft \trsl{Marek is eating bread for the child's benefit.}//
\endgl
\a
\begingl
\gla Byl Marcí travě pjaštébime.//
\glb child.\Dir{} Marek.\Ind{} bread.\Ind{} eat-\Appl{}-\Prog{}//
\glft \trsl{Marek is eating bread for the child's benefit.}//
\endgl
\xe

The applied child alone bears direct case. Marek and the bread remain arguments
of the eating event, although their shared indirect case does not distinguish
actor from theme; the predicate and context supply those relations. Neither
clause means that the child eats. Causative \ird{pjáštrime}, by contrast,
can describe Marek's bringing about the child's eating
(§\,\ref{sec:aspect-caus}). A conventional \trsl{feed} reading is not part of
the productive semantics of \ird{-éb-}.

Two independently direct arguments are excluded from this simple applicative
frame:

\ex
\begingl
\gla *Marek byl pjaštébime.//
\glb Marek.\Dir{} child.\Dir{} eat-\Appl{}-\Prog{}//
\glft \trsl{(intended: Marek is eating for the child's benefit.)}//
\endgl
\xe

An intransitive activity permits the same benefactive derivation without a
lexical theme. The following constructed examples illustrate that difference
and the ERG-domain frame of the transitive applicative:

\pex\label{ex:val-appl-ergative}
\a
\begingl
\gla Byl Marcí zemacébime.//
\glb child.\Dir{} Marek.\Ind{} run-\Appl{}-\Prog{}//
\glft \trsl{Marek is running for the child's benefit.}//
\endgl
\a
\begingl
\gla Byl Marku travě pjaštébik.//
\glb child.\Dir{} Marek.\Trs{} bread.\Ind{} eat-\Appl{}-\Pf{}//
\glft \trsl{Marek ate bread for the child's benefit.}//
\endgl
\xe

\ird{-éb-} also admits role interpretations constrained by the predicate:
a recipient requires transfer or availability of a theme, a goal requires
motion or placement, and an instrument requires an activity in which that
participant serves as an implement. Locative uses require a lexically or
constructionally licensed spatial relation; the location of an event and its
destination are distinct relations. Discourse salience can motivate selection
of such a construction but cannot supply an otherwise unlicensed semantic
role. These restrictions distinguish Iridian's interpretations within the
broader typological variation in applicative bases and roles
\cite[§§2, 4]{polinsky2013appl}.

The goal and instrumental extensions can be illustrated with controlled
constructed pairs. Directed-motion \ird{stoj-} combines with a goal phrase
in the first example below; applicativisation makes that same town the
direct-case participant in the second. Both retain contemplative aspect:

\pex\label{ex:val-appl-goal}
\a
\begingl
\gla Marek vruham de stojáž.//
\glb Marek.\Dir{} town.\Ind{} to go-\Ctpl{}//
\glft \trsl{Marek intends to go to the town.}//
\endgl
\a
\begingl
\gla Vruh Marcí stojébaž.//
\glb town.\Dir{} Marek.\Ind{} go-\Appl{}-\Ctpl{}//
\glft \trsl{Marek intends to go to the town.}//
\endgl
\xe

The town is the destination of projected motion, without an implication that
it benefits from that motion. Instrumental \ird{se}, already illustrated
with the antipassive, instead relates an implement to an activity. In the
following constructed pair \ird{rán}/\ird{ráne} supplies a controlled
case alternation (§\,\ref{sec:nm-stem-grades}); the context identifies the
hand as the means by which Marek eats:

\pex\label{ex:val-appl-instrument}
\a
\begingl
\gla Marek travě ráne se pjaštíme.//
\glb Marek.\Dir{} bread.\Ind{} hand.\Ind{} \textsc{inst} eat-\Prog{}//
\glft \trsl{Marek is eating bread by hand.}//
\endgl
\a
\begingl
\gla Rán Marcí travě pjaštébime.//
\glb hand.\Dir{} Marek.\Ind{} bread.\Ind{} eat-\Appl{}-\Prog{}//
\glft \trsl{Marek is eating bread by hand.}//
\endgl
\xe

Neither extension changes the lexical event: the hand does not become an
eater, and the town does not become a moving participant. The instrumental
interpretation depends on an implement-compatible event, just as the goal
interpretation depends on directed motion; unrestricted exchange of these
roles is excluded by those selectional conditions. Recipient and locative
extensions remain restricted to compatible predicates, whose full lexical
distribution is unspecified.

Directive Mood selects the lexical theme where one is retained, so that an
applied beneficiary receives indirect case (§\,\ref{sec:mood-directive-valency}):

\pex
\a
\begingl
\gla Trava bylam pjaštébím.//
\glb bread.\Dir{} child.\Ind{} eat-\Appl{}-\Imp{}//
\glft \trsl{Eat bread for the child's benefit.}//
\endgl
\a
\begingl
\gla Trava bylam pjaštébká.//
\glb bread.\Dir{} child.\Ind{} eat-\Appl{}-\Hort{}//
\glft \trsl{Let's eat bread for the child's benefit.}//
\endgl
\xe

The beneficiary is required overtly in these directive constructions. Its
indirect case reflects the higher Mood-conditioned selection rule, although
Appl continues to introduce it as an argument.


The productive benefactive is analysed as a high applicative, whose added
participant is related to an event. Pylkkänen distinguishes this relation from
low applicatives, which relate individuals through transfer of possession
\cite[ch.~2]{pylkkanen2008}. Iridian's benefactive introduces no possession
entailment, and the grammar assigns its productive interpretation to a single
Appl head. Recipient readings are predicate-conditioned extensions rather than
evidence for an independently productive low-applicative position.

A schematic interpretation is
\[
 \mathrm{Appl}_{\mathrm{BEN}}
 =\lambda P.\lambda y.\lambda e\,
 [P(e)\land\operatorname{Benefit}(e,y)],
\]
where $P$ is an event predicate and $y$ is the beneficiary. Appl combines with
the event predicate below the layer introducing its actor; both lie below an
optional Caus head and the outer orientation head. The constituent order is
head-final throughout:

\begin{center}
\begin{forest}
for tree={align=center}
[\synhead{ApplP}
  [{\synhead{DP}\\{\footnotesize beneficiary}}]
  [{Appl$'$}
    [\synhead{VP} [{\ldots}]]
    [{Appl\\\ird{-éb-}}]
  ]
]
\end{forest}
\end{center}

Appl introduces the beneficiary; the completed clause's licensing configuration
determines case and pivot privileges. A simple indicative applicative selects
the beneficiary in either alignment domain, whereas higher Caus or directive
Mood can supersede that selection. The outer orientation head is zero in a simple applicative, where its
selection follows Appl; \ird{-éb-} selects perfective \ird{-ik} and
licenses the ERG-domain frame without an additional \ird{-n-}. A higher
causative creates its own orientation domain and requires overt Voice
\ird{-an-} in the ERG frame. Iridian excludes a
second applicative and the combination of applicative and antipassive within
the same predicate.

An applicative makes an event-related participant available for pivot-restricted
constructions, including participant nominalisation and relative formation.
Continued reference to a beneficiary may favour this choice, although a topic
relation is neither required nor supplied by the suffix. The beneficiary
interpretation comes from the applicative construction; discourse salience
does not by itself license an instrumental, goal or stimulus reading.

\subsection{The passive}\label{sec:vm-passive}

The NOM-domain passive combines patient-selecting \ird{-n-}
(\ird{-an-} after a rhotic causative) with continuous
\ird{-a}, progressive \ird{-ime}, or contemplative \ird{-áž}/\ird{-až}, selected by the local
quantity specification. Its patient
receives direct case, while an expressed actor receives indirect case as an
agentive dependent licensed by the passive construction. An unexpressed actor
can be contextually recovered or left unspecified. The aspectual domain remains
NOM, although active actor selection is overridden.

The indirect agent differs in licensing from the ordinary indirect patient of
an active NOM-domain transitive: the latter is an internal argument of the
lexical predicate, whereas the former is an optional external dependent of the
passive configuration. Their shared case ending does not establish identical
argument status. Further diagnostics of the passive agent's binding and control
properties are not established by the case alternation alone.

\pex
\a
\begingl
\gla Travó Marcí pjaštníme.//
\glb bread.\Dir{}.\Def{} Marek.\Ind{} eat-\textsc{pass}-\Prog{}//
\glft \trsl{The bread is being eaten by Marek.}//
\endgl
\a
\begingl
\gla Travó pjaštníme.//
\glb bread.\Dir{}.\Def{} eat-\textsc{pass}-\Prog{}//
\glft \trsl{The bread is being eaten.}//
\endgl
\xe

\ird{Marcí} is the indirect form of \ird{Marek}; transitive \ird{Marku}
belongs to the ERG-domain active frame. The same outer head follows the
causative in the constructed passive below, where Marek is the optional causer
dependent and the child remains the embedded actor:

\ex
\begingl
\gla Trava Marcí bylam pjášt-r-an-ime.//
\glb bread.\Dir{} Marek.\Ind{} child.\Ind{} eat-\Caus{}-\textsc{pass}-\Prog{}//
\glft \trsl{Bread is being caused by Marek to be eaten by the child.}//
\endgl
\xe

The passive privileges the bread for
participant nominalisation and direct-gap relativisation. It can thereby
support continued reference to that participant, although aboutness topicality
requires its independent discourse conditions.


\section{Causatives and outer Voice}\label{sec:aspect-caus}

The rhotic causative adds a causer while retaining the base's independently
licensed participants. Compact \ird{-r-} and expanded \ird{-ér-}/\ird{-er-}
differ in their local distribution of strength, while a lexically restricted
expanded class also selects front-grade PAL. These valency increases have a
comparison in Dixon's morphological causative type \cite{dixon2000};
the allomorphs, quantities, and case rules below specify Iridian.

\subsection{Rhotic causative morphology}
\label{sec:causative-grades}

Allomorph selection precedes quantity assignment. Compact \ird{-r-}
selects G1 on an eligible local stem nucleus, giving \ird{pjáštr-},
\ird{vádr-}, and \ird{zemácr-}. It does not change the preceding
thematic consonant: \ird{pjáštrime} retains /t/, and
\ird{zemácrime} retains /ts/. The compact onset clusters /tr/ and /br/
follow the ordinary syllabification rules
(§\,\ref{sec:cluster-repair}).

Expanded \ird{-er-} occurs where the compact sequence cannot be
syllabified and in specified lexical classes, including \ird{ščen-},
stative stems in \ird{-n-}, \ird{stoj-}, \ird{nočk-}, and the
replacement antipassive \ird{pjášč-}. With a short input stem, the
expanded exponent carries strength as \ird{-ér-}:
\ird{ščenérime}, \ird{svornérime}. Already strong \ird{pjášč-}
instead takes short \ird{-er-}, as in \ird{pjáščeranik}.
The listed compact alternative \ird{ščénrime} retains G1 on the root;
its existence does not make every expanded-class stem freely compact.

Before consonant-initial directive Mood, the expanded allomorph likewise
supplies the causative vowel: \ird{pjaštérká}, \ird{pjaštérku}.
These forms select expanded Caus before assigning strength to it.
Compact \ird{pjáštrím} and \ird{pjáštré}, by contrast, retain
their root G1 alongside the independently long Mood vowel. The distribution
does not follow a ban on two long vowels in a word.

The mutating expanded class is lexical. Its front-grade consonantal
specification is independent of the surface vowel quantity:
\ird{vad-} with that causative gives \ird{vaďérime}, beside compact
\ird{vádrime}. PAL supplies /ɟ/, written \ird{ď}; G1 belongs to
the suffix vowel, so the root remains short. Long \ird{é} replaces
the short-vowel spelling \ird{ě} in the strong expanded exponent.
The front and plain expanded classes can therefore share the spelling
\ird{-ér-}/\ird{-er-} while selecting different consonantal
correspondences. No spelling-based rule freely assigns PAL to another
causative class or changes /r/ to learned \ird{ř}.

An intervening valency head establishes a separate domain. Protected
applicative \ird{-éb-} retains its quantity in \ird{pjaštébrime}
and prevents outer Caus from lengthening the original root. The same
locality prevents a causative outside antipassive morphology from applying
PAL or assigning a new root grade through that inner head.
Morphology-sensitive palatalisation has a comparative parallel in the Polish
suffix classes discussed by Rubach \cite[§§2--3]{rubach2019}; the lexical
classes and local quantitative specifications stated here belong to Iridian.

Outer patient-selecting Voice takes \ird{-an-} immediately after any
rhotic causative, wherever that orientation head is overt, and ordinary
\ird{-n-} elsewhere. Thus \ird{pjášt-r-an-ik} yields
\ird{pjáštranik}, while \ird{pjášt-r-an-ime} yields passive
\ird{pjáštranime}. The short aspect endings preserve Caus-selected
strength. Expanded \ird{pjášč-er-an-ik} gives \ird{pjáščeranik},
and front-grade expanded \ird{vaď-ér-an-ik} gives \ird{vaďéranik}.
The causative and orientation remain distinct exponents, with no metathesis
or joint causative--Voice morpheme.

In the Distributed Morphology analysis, these causative and orientation
allomorphs are Vocabulary Items with lexical or morphological conditioning,
following Embick and Noyer's account of contextual allomorphy
\cite[§2.3]{embicknoyer2007}. Their lexical membership does not follow from
a root's meaning.

\begin{table}[htbp]
\footnotesize
\caption{Constructed causative forms of compact-class \ird{pjašt-}.}
\label{tab:rhotic-causative-paradigm}
\begin{tabularx}{\textwidth}{L L L L}
\toprule
Category & Form & Category & Form \\
\midrule
Continuous & \ird{pjáštra} & Progressive & \ird{pjáštrime} \\
Contemplative & \ird{pjáštraž} & Perfective & \ird{pjáštranik} \\
Retrospective & \ird{pjáštranav} & Imperative & \ird{pjáštrím} \\
Prohibitive & \ird{pjáštré} & Hortative & \ird{pjaštérká} \\
PFV imperative & \ird{pjáštrovím} & Sequential & \ird{pjáštru}, \ird{pjáštrení} \\
\bottomrule
\end{tabularx}
\end{table}

\subsection{Causative argument frames}

In an active NOM-domain causative, the causer is direct and all other arguments
are indirect. In an ERG-domain causative, the causer is transitive and one
embedded participant is direct: the retained lexical theme takes precedence,
followed by an applied participant if there is no lexical theme, and otherwise
the causee. Any remaining causee or applied participant is indirect. This
hierarchy belongs to the causative frame, so an embedded applicative's simple
beneficiary selection does not independently assign a second direct case.
Here \term{retained theme} designates the separate patient of a transitive
base. An intransitive's sole participant becomes the causee even when its
lexical role is theme-like, as with \ird{ščen-} \trsl{arrive}; that role
does not create an additional patient ahead of an applied participant in the
causative hierarchy.

\begin{table}[htbp]
\footnotesize
\caption{Active indicative causative frames. A dash indicates an absent role.}
\label{tab:causative-frames}
\begin{tabularx}{\textwidth}{L C C C C}
\toprule
Base and domain & Causer & Causee & Theme & Applied \\
\midrule
Intransitive, NOM & \textsc{dir} & \textsc{ind} & --- & --- \\
Intransitive, ERG & \textsc{trs} & \textsc{dir} & --- & --- \\
Transitive, NOM & \textsc{dir} & \textsc{ind} & \textsc{ind} & --- \\
Transitive, ERG & \textsc{trs} & \textsc{ind} & \textsc{dir} & --- \\
Applied transitive, NOM & \textsc{dir} & \textsc{ind} & \textsc{ind} & \textsc{ind} \\
Applied transitive, ERG & \textsc{trs} & \textsc{ind} & \textsc{dir} & \textsc{ind} \\
Applied intransitive, NOM & \textsc{dir} & \textsc{ind} & --- & \textsc{ind} \\
Applied intransitive, ERG & \textsc{trs} & \textsc{ind} & --- & \textsc{dir} \\
Antipassive, ERG & \textsc{trs} & \textsc{dir} & --- & --- \\
\bottomrule
\end{tabularx}
\end{table}

The base contrasts can be compared directly in the following constructed
paradigm. \ird{Zemac-} denotes an actor's activity, whereas \ird{ščen-}
denotes its sole participant's arrival; their ordinary case frames coincide.
The state \ird{svorn-} independently selects perfective \ird{-ik};
the eventive inchoative \ird{svornél-} instead selects free \ird{-ek}.
The transitive base alone requires \ird{-n-} in the ERG frame.

\begin{table}[htbp]
\footnotesize
\caption{Constructed ordinary base clauses in the two alignment domains.}
\label{tab:valency-base-contrasts}
\begin{tabularx}{\textwidth}{l L L}
\toprule
Base & NOM: progressive & ERG: perfective \\
\midrule
\ird{pjašt-} & \ird{Marek travě pjaštíme.} & \ird{Trava Marku pjaštník.} \\
\ird{zemac-} & \ird{Marek zemácime.} & \ird{Marek zemácek.} \\
\ird{ščen-} & \ird{Byl ščénime.} & \ird{Byl ščének.} \\
\ird{svornél-} & \ird{Problem svornélime.} & \ird{Problem svornélek.} \\
\bottomrule
\end{tabularx}
\end{table}

In the transitive row the agent is direct in NOM and transitive in ERG,
while the theme is indirect in NOM and direct in ERG. Each intransitive row
has one core participant, which is direct and privileged in both domains.
The progressive inchoative describes developing scalar change, while its
perfective bounds that change without independently fixing a threshold; stative causativisation adds
a distinct causer (§\,\ref{sec:stative-valency}).

Complex forms below retain visible morphological boundaries where these
clarify suffix order; the hyphens are analytical notation. The following
constructed pairs make the causee and any retained theme overt:

\pex\label{ex:val-caus-intransitive}
\a
\begingl
\gla Marek bylam zemácrime.//
\glb Marek.\Dir{} child.\Ind{} run-\Caus{}-\Prog{}//
\glft \trsl{Marek is making the child run.}//
\endgl
\a
\begingl
\gla Byl Marku zemác-r-an-ik.//
\glb child.\Dir{} Marek.\Trs{} run-\Caus{}-\textsc{trs}-\Pf{}//
\glft \trsl{Marek made the child run.}//
\endgl
\xe

The activity predicate \ird{zemac-} \trsl{run} supplies an actor causee.
By comparison, the sole participant of \ird{ščen-} \trsl{arrive} undergoes a
change of location. The grammar treats the former as unergative and the latter
as unaccusative, using that distinction to characterise the base relation;
both follow the same causative case rule. The following constructed paradigm
includes the underived ERG-domain intransitive, which needs no transitive
\ird{-n-}, and causatives in both alignment domains:

\pex\label{ex:val-unaccusative}
\a
\begingl
\gla Byl ščének.//
\glb child.\Dir{} arrive-\Pf{}//
\glft \trsl{The child arrived.}//
\endgl
\a
\begingl
\gla Marek bylam ščenérime.//
\glb Marek.\Dir{} child.\Ind{} arrive-\Caus{}-\Prog{}//
\glft \trsl{Marek is causing the child to arrive.}//
\endgl
\a
\begingl
\gla Byl Marku ščen-ér-an-ik.//
\glb child.\Dir{} Marek.\Trs{} arrive-\Caus{}-\textsc{trs}-\Pf{}//
\glft \trsl{Marek caused the child to arrive.}//
\endgl
\xe

\pex\label{ex:val-caus-transitive}
\a
\begingl
\gla Marek bylam travě pjáštrime.//
\glb Marek.\Dir{} child.\Ind{} bread.\Ind{} eat-\Caus{}-\Prog{}//
\glft \trsl{Marek is making the child eat bread.}//
\endgl
\a
\begingl
\gla Travó Marku bylam pjášt-r-an-ik.//
\glb bread.\Dir{}.\Def{} Marek.\Trs{} child.\Ind{} eat-\Caus{}-\textsc{trs}-\Pf{}//
\glft \trsl{Marek made the child eat the bread.}//
\endgl
\xe

The rhotic causative expresses causation without specifying coercion, permission, or the
causer's degree of intervention. The causee's agency, the lexical event, and
context support more particular readings; viewpoint aspect does not entail
causal directness or control. Thus the perfective transitive example can
receive a permissive interpretation in a context where Marek enables the
child to eat, while retrospective \ird{pjášt-r-an-av} can describe direct
compulsion when that relation is supplied by context. Both preserve the same
causative argument frame.

An applicative precedes the causative. The following constructed clauses use
plural \ird{býlac} for the causee and singular \ird{bylam} for the beneficiary,
so that all four participants are represented:

\pex\label{ex:val-caus-appl}
\a
\begingl
\gla Marek býlac bylam travě pjaštébrime.//
\glb Marek.\Dir{} child.\Ind{}.\Pl{} child.\Ind{} bread.\Ind{} eat-\Appl{}-\Caus{}-\Prog{}//
\glft \trsl{Marek is making the children eat bread for the child's benefit.}//
\endgl
\a
\begingl
\gla Trava Marku býlac bylam pjaštéb-r-an-ik.//
\glb bread.\Dir{} Marek.\Trs{} child.\Ind{}.\Pl{} child.\Ind{} eat-\Appl{}-\Caus{}-\textsc{trs}-\Pf{}//
\glft \trsl{Marek made the children eat bread for the child's benefit.}//
\endgl
\xe

Their scope is $\mathrm{CAUS}[\mathrm{APPL}[\mathrm{EAT}]]$: the child benefits
from the embedded eating event. A beneficiary of the causing event itself can
be expressed separately with \ird{my}; the suffix order does not license an
outer applicative reading. Multiple indirect-case arguments need not occur in
all uses, but an omitted causee requires the ordinary contextual licensing of
argument omission. The morphology itself does not introduce an arbitrary
silent \trsl{someone}.

A causative of an applied intransitive lacks a separate lexical patient.
The following constructed pair therefore selects the causer in the
progressive and the beneficiary in the perfective, leaving the running
children in indirect case in both:

\pex\label{ex:val-caus-appl-intransitive}
\a
\begingl
\gla Marek býlac bylam zemacébrime.//
\glb Marek.\Dir{} child.\Ind{}.\Pl{} child.\Ind{} run-\Appl{}-\Caus{}-\Prog{}//
\glft \trsl{Marek is making the children run for the child's benefit.}//
\endgl
\a
\begingl
\gla Byl Marku býlac zemacéb-r-an-ik.//
\glb child.\Dir{} Marek.\Trs{} child.\Ind{}.\Pl{} run-\Appl{}-\Caus{}-\textsc{trs}-\Pf{}//
\glft \trsl{Marek made the children run for the child's benefit.}//
\endgl
\xe

The perfective child is the beneficiary, while plural \ird{býlac}
expresses the causee. The contrast with the preceding transitive pair
follows from the presence or absence of a retained patient, rather than from
a change in the benefactive relation contributed by Appl.

An antipassive base retains its actor as the causee while its lexical patient
remains unavailable. Causativisation therefore produces a new two-participant
predicate without reversing the lower reduction:

\ex\label{ex:val-caus-antipassive}%
\begingl
\gla Byl Marku pjáščeranik.//
\glb child.\Dir{} Marek.\Trs{} eat-\Antip{}-\Caus{}-\textsc{trs}-\Pf{}//
\glft \trsl{Marek made the child eat.}//
\endgl
\xe

The ERG-domain licensing restriction on antipassive remains operative in this
configuration; a NOM-domain counterpart containing \ird{-aš-} is excluded.
The direct-case child is the causee, never a restored lexical patient.

\subsection{Causative structure and participant selection}
\label{sec:vm-caus}

Caus combines with an already formed event predicate, introducing its causer
above the base actor and any applied argument. The outer orientation head
subsequently selects the derived predicate's privileged participant, while
Asp.Align supplies the alignment domain. The distinction between productive
causative structure and other verbal heads follows the analytical issues
surveyed by Harley \cite[§1.4]{harley2017}; assigning causer introduction
to Caus and patient selection to higher Voice is the Iridian analysis.
In the transitive causative \ird{pjáštranik}, the lexical patient originates
inside VP, the lower actor layer introduces the causee, and Caus introduces
the causer above that event predicate. The following head-final tree displays
those thematic origins; ordinary finite placement subsequently determines
where the arguments are pronounced (§\,\ref{sec:vm-pivot-derivations}). The
verbal categoriser and higher inflectional layers are omitted.

\begin{center}
\begin{forest}
for tree={align=center}
[\synhead{AspP}
  [\synhead{VoiceP}$_{\mathrm{orient}}$
    [\synhead{CausP}
      [{\synhead{DP}\\{\footnotesize causer}}]
      [{Caus$'$}
        [{inner actor layer}
          [{\synhead{DP}\\{\footnotesize actor/causee}}]
          [{actor$'$}
            [\synhead{VP}
              [{\synhead{DP}\\{\footnotesize patient}}]
              [{V\\\ird{pjášt-}}]
            ]
            [{actor\\$\varnothing$}]
          ]
        ]
        [{Caus\\\ird{-r-}}]
      ]
    ]
    [{Voice\\\ird{-an-}}]
  ]
  [{Asp.Align\\\ird{-ik}}]
]
\end{forest}
\end{center}

The corresponding morphological constituency is
\[
 \bigl[\bigl[\bigl[\ird{pjášt}-\textsc{caus}\bigr]
 -\textsc{orient}\bigr]-\textsc{pfv}\bigr]
 \quad\longmapsto\quad \ird{pjášt-r-an-ik}.
\]
This mapping assumes successive suffixal attachment within one verbal word;
head-final complement order alone would not supply that attachment operation.
The silent categoriser and actor head have no segment in the displayed form.
Compact Caus selects the strong root grade \ird{pjášt-}, its rhotic exponent
conditions orientation \ird{-an-}, and the derived stem selects short
perfective \ird{-ik} (§\,\ref{sec:causative-grades}). In the completed
ERG frame, orientation selects the retained patient as direct, the completed
aspect--orientation configuration licenses the causer as transitive, and the
causee is indirect, as illustrated
by \ird{Travó Marku bylam pjášt-r-an-ik}
(\ref{ex:val-caus-transitive}). These case relations follow the construction's
licensing rules rather than the spelling of its suffixes.

The overt outer \ird{-an-} licenses the selected embedded participant in an
ERG causative. In the NOM domain, zero outer orientation gives the active
causer-oriented frame; overt \ird{-an-} selects the embedded theme, then an
applied participant if no theme is retained, and otherwise the causee. It
licenses the causer as an optional indirect agentive dependent. This is the passive of the
derived causative. An embedded causee remains an argument of its own predicate;
passivisation of the higher causing event does not suppress that relation.

The established contrasts support a set of separable checks on this mapping.
Table~\ref{tab:vm-mirror-diagnostics} compares forms already specified above;
the rightmost column records the structural or local morphological relation
which an alternative derivation must preserve.

\begin{table}[htbp]
\footnotesize
\caption{Diagnostics of verbal composition, using established constructed forms.}
\label{tab:vm-mirror-diagnostics}
\begin{tabularx}{\textwidth}{L l L}
\toprule
Construction & Form & Relation to preserve \\
\midrule
Simple ERG transitive & \ird{pjaštník}
 & Patient direct, actor transitive; Class B quantity on aspect \\
Active NOM causative & \ird{pjáštrime}
 & Causer direct; causee and retained patient indirect; Caus selects root G1 \\
ERG causative & \ird{pjáštranik}
 & Retained patient direct; causer transitive; causee indirect; rhotic Caus conditions \ird{-an-} \\
ERG applied causative & \ird{pjaštébranik}
 & Retained patient direct, causer transitive; causee and beneficiary indirect; Appl below Caus protects the root from outer G1 \\
\bottomrule
\end{tabularx}
\end{table}

\noindent The applicative--causative contrast additionally fixes the
benefactive relation inside the caused event: \ird{pjaštébranik} has
$\mathrm{CAUS}[\mathrm{APPL}[\mathrm{EAT}]]$ scope, without an outer
benefactive interpretation (\ref{ex:val-caus-appl}). Applying indicative
negation to each row predicts immediate left-edge \ird{za-}, with the listed
participant inventory, suffix order and local grades preserved under negation.
These negative complex counterparts are predictions of the composition and
realisation rules, rather than additional independently established paradigms;
the simple \ird{zapjaštník} already illustrates the prefixal mapping
(§\,\ref{sec:neg-allomorphy}). Their interpretation negates the complete
aspectual causative or applicative proposition, without giving \ird{za-} a
position between the inner derivational heads.

\subsection{Participant dependencies and causative locality}
\label{sec:vm-pivot-derivations}

The structures below separate thematic origin from the selection and case
requirements in Table~\ref{tab:pivot-case-obligations}. They are schematic
argument structures, not surface-order trees. $\operatorname{Act}(A,Q)$
abbreviates the inner actor-introducing configuration over predicate $Q$;
$\operatorname{Caus}(C,Q)$ abbreviates Caus with its causer. These are
expository abbreviations for the argument layers already described, not new
functional projections. In all five structures the theme originates as V's
complement; selecting it does not alter that origin.

\begin{gather*}
\text{ERG transitive:}\quad
 \operatorname{Orient}_{T}[\operatorname{Act}(A,[_{\mathrm{VP}}T\ V])]
 \\ \Rightarrow T_{\mathrm{DIR}}\ A_{\mathrm{TRS}}\ V_{\mathrm{PFV}};\\[3pt]
\text{active NOM:}\quad
 \operatorname{Orient}_{A}[\operatorname{Act}(A,[_{\mathrm{VP}}T\ V])]
 \\ \Rightarrow A_{\mathrm{DIR}}\ T_{\mathrm{IND}}\ V_{\mathrm{IPFV}};\\[3pt]
\text{simple applicative:}\quad
 \operatorname{Orient}_{B}[\operatorname{Act}(A,[_{\mathrm{ApplP}}B\ [_{\mathrm{VP}}T\ V]\ \mathrm{Appl}])]
 \\ \Rightarrow B_{\mathrm{DIR}}\ A_{\mathrm{IND/TRS}}\ T_{\mathrm{IND}}\ V;\\[3pt]
\text{ERG causative:}\quad
 \operatorname{Orient}_{T}[\operatorname{Caus}(C,\operatorname{Act}(E,[_{\mathrm{VP}}T\ V]))]
 \\ \Rightarrow T_{\mathrm{DIR}}\ C_{\mathrm{TRS}}\ E_{\mathrm{IND}}\ V;\\[3pt]
\text{ERG applied causative:}\\
 \operatorname{Orient}_{T}[\operatorname{Caus}(C,\operatorname{Act}(E,[_{\mathrm{ApplP}}B\ [_{\mathrm{VP}}T\ V]\ \mathrm{Appl}]))]
 \\ \Rightarrow T_{\mathrm{DIR}}\ C_{\mathrm{TRS}}\ E_{\mathrm{IND}}\ B_{\mathrm{IND}}\ V.
\end{gather*}

Here the right-hand sequences record the ordinary patterns illustrated in the
existing paradigms, rather than deriving their order by fiat. The third row
uses indirect actor case in NOM and transitive actor case in ERG. Higher Caus
changes selection in the fifth row: $T$ retains priority over $B$, although
$B$ lies higher inside the embedded predicate. With no retained $T$, the
applied causative selects $B$; with neither $T$ nor $B$, including an
antipassive base, it selects $E$. In a NOM passive of the causative, the
selected embedded participant is direct, the causer is an optional indirect
dependent, and the causee retains its embedded argument relation. Directive
Mood preserves a retained theme as direct in both reduced aspectual domains.
These extensions preserve the argument inventory while changing the completed
construction's licensing configuration.

Ordinary finite placement derives these selected-first sequences in
 \synhead{OrdP} (§\,\ref{sec:pivot-placement}). Its first search targets the
 participant selected by the completed orientation--Asp.Align--Mood
 configuration, rather than the nearest DP without restriction; subsequent
 searches place the remaining eligible constituents in their lower order.
 Thus an ergative transitive places $T$ before $A$, a simple applicative
 places $B$ before $A$ and $T$, and an ergative transitive causative places
 $T$ before $C$ and $E$. Tucking each later goal beneath the earlier ordinary
 edge occupants preserves the residual order. The selection dependency
 follows the licensing configurations in
 Table~\ref{tab:pivot-case-obligations}; Ord does not supply an unexplained
 participant-selection feature or recompute case from the resulting string.
 The theme's lower occurrence remains available for binding and interpretation.

Nominal intervention and phase access remain distinct conditions. The first
 Ord search excludes a closer actor or causer because that participant is not
 selected by the completed licensing configuration. Phase access instead
 depends on which lower heads delimit transfer domains. The finite placement
 analysis assigns an edge requirement to each actor- or causer-introducing
 head that is treated as a phase within the same finite clause. Before its
 complement is transferred, this requirement registers the eligible ordinary
 constituents at an ordered accessible edge: an eligible constituent already
 at the edge remains accessible, while lower goals move there nearest first,
 tucking beneath the earlier edge occupants. Searches access the highest
 accessible occurrence of each chain and ignore its already registered lower
 copies. The ordering of internally and externally merged specifiers is a
 substantive choice: Paillé's analysis of pivot-third order argues for the
 possibility of tucking below an externally merged actor
 \cite[§5.1]{paille2021}. The Iridian rule preserves that actor's
 precedence at the intermediate edge; final pivot priority is established
 separately by Ord. This comparison concerns specifier order; the
 equidistance-based binding analysis in that account is not adopted for the
 Iridian ordinary edge. The preliminary steps preserve the original residual
 order and do not select the pivot, whose identity can
 still be affected by higher orientation or Mood. Ord subsequently attracts
 the participant selected by the completed configuration before placing the
 residual constituents. A non-phasal head introduces no additional escape
 step. For an ergative transitive causative, the selected theme's path is:
\[
 T\ \longrightarrow\ \mathrm{edge}(\text{inner actor domain})
 \ \longrightarrow\ \mathrm{edge}(\text{causer domain})
 \ \longrightarrow\ \mathrm{OrdP},
\]
with each intermediate step present only if that domain is a phase.

This edge requirement is a property of the movement analysis; the surface
 paradigms do not independently establish which argument-building heads are
 phases. It licenses the stated clause-local routes without treating every
 complement as transparent. Predicate heads and operator particles are not
 ordinary placement goals, and internal constituents of a selected finite
 clause, converb or nominal do not enter the matrix search. Moving the whole
 selected clause preserves that boundary. The explicit finite placement
 derivations and the lower chain conditions are given in
 §§\,\ref{sec:pivot-placement} and~\ref{sec:pivot-scope-connectivity}.

\subsection{Outer Voice and patient orientation}
\label{sec:vm-n}

\ird{-n-} realises outer patient-selecting Voice, following causative
\ird{-r-}, \ird{-er-}, or \ird{-ěr-} wherever both occur,
with the conditioned Voice allomorph \ird{-an-}. With a simple transitive base it
selects the
lexical patient; with a causative it selects according to the hierarchy in
Table~\ref{tab:causative-frames}. Its interpretation as patient selection thus
extends to a causee that is the affected participant of an intransitive-base
causative. This higher orientation head is distinct from inner antipassive
Voice and from argument-introducing Appl and Caus.

In the ERG-aspect domain, \ird{-n-} licenses the ordinary two-participant
transitive frame: the selected participant is direct and the highest actor is
transitive. Simple applicatives have their own overt licensing through
\ird{-éb-}; simple antipassives instead retain only an actor, which receives
direct case. The contrast between overt \ird{-n-} and a bare perfective on a
transitive base is therefore a restriction on the completed licensing frame,
not evidence that patient argumenthood universally requires an overt suffix.

In the NOM-aspect domain, \ird{-n-} overrides active actor selection and
licenses a different external-argument configuration: the patient is direct,
while the highest actor can occur as an optional indirect agentive dependent.
The grammar calls this configuration \term{passive}. Patient selection and
external-argument licensing are separate parts of its specification; assigning
direct case to a patient alone would not establish the latter property.
The relative height of Voice accounts for \ird{pjášt-r-an-ik} without
metathesis of Voice and Caus.

A causative introduces a causing event and a causer, while the lexical predicate
and context determine whether the causal relation involves inducement,
permission, delegated agency or physical intervention. Perfective and
retrospective viewpoints can present the same relation from different temporal
perspectives without encoding different degrees of force. Directive Mood
likewise changes the construction's orientation without changing the causal
relation (§\,\ref{sec:mood-directive-valency}).
